% GENERATED FILE. Edit canonical lessons, then run npm run generate:books.
% canonical-source-hash: fnv1a64:9364ca7633f1017a
% canonical-lessons: IT-C382-ripassare, IT-C382-interrogazione, IT-C382-memorizzare, IT-C382-appunti, IT-C382-manuale

\chapter{Prima dell'esame}
\label{ch:it-prima-dell-esame}
\hlchaptermodality{382}

\begin{chapteropening}
\textbf{By the end of this chapter:} \emph{I can talk about getting ready for a test: revising, an oral test, learning by heart, my notes and the textbook.}

Complete the last lesson of chapter 382: I can talk about getting ready for a test: revising, an oral test, learning by heart, my notes and the textbook.
\end{chapteropening}

\section[ripassare]{ripassare — to revise, to go over (again)}

\label{lesson:IT-C382-ripassare}



\begin{quote}
Before the new one: say the Italian for the doctorate, then the Italian for the exam session.
\end{quote}

\subsection*{You'll want to know: ripassare}
\textbf{ripassare} — \textquotedblleft{}to revise, to go over (again)\textquotedblright{}.

\begin{quote}
Stasera ripassiamo la storia per l'esame.
\end{quote}

\begin{quote}
\emph{Tonight we are going over history for the exam.}
\end{quote}

It also means to come by again: \emph{ripassa domani}.

\begin{cousinweb}
\emph{ripassare} is built from \emph{ri-} \textquotedblleft{}again\textquotedblright{} and \emph{passare} \textquotedblleft{}to pass\textquotedblright{}, a verb inherited from Vulgar Latin \emph{passāre}, built on Latin \emph{passus} \textquotedblleft{}a step\textquotedblright{} (Lewis \& Short).
\end{cousinweb}

\begin{grammarlens}[title={What you've built}]
One more word for telling, explaining and describing at B1.
\end{grammarlens}

\subsection*{Your turn}
\begin{itemize}
\raggedright
  \item \emph{Say it:} \emph{ripassare}
  \item \emph{Say it:} \emph{ripassare}, once more
  \item \emph{Say it:} a sentence of your own with \emph{ripassare}
  \item \emph{Recall:} say the Italian for the doctorate, then the Italian for the exam session, then say \emph{ripassare} again
\end{itemize}

\begin{tcolorbox}[breakable,colback=teal!4,colframe=teal!35!black,title={Before you move on}]
What does ripassare mean? (\textquotedblleft{}To revise, to go over (again)\textquotedblright{}.) Say it once more.
\end{tcolorbox}
\section[l'interrogazione]{l'interrogazione — the oral test (at school); the questioning}

\label{lesson:IT-C382-interrogazione}



\begin{quote}
Before the new one: say the Italian for the exam session, then the Italian for to revise.
\end{quote}

\subsection*{You'll want to know: l'interrogazione}
\textbf{l'interrogazione} — \textquotedblleft{}the oral test (at school); the questioning\textquotedblright{}.

\begin{quote}
Se domani c'è l'interrogazione di storia, stasera studio.
\end{quote}

\begin{quote}
\emph{If there is an oral history test tomorrow, I will study tonight.}
\end{quote}

At school the teacher questions one pupil aloud in front of the class: \emph{essere interrogato}.

\begin{cousinweb}
\emph{interrogazione} is a learned borrowing from Latin \emph{interrogātiō} \textquotedblleft{}a questioning, an examination\textquotedblright{} (EtymDB), from \emph{interrogāre} \textquotedblleft{}to question\textquotedblright{}, built from \emph{inter} \textquotedblleft{}between\textquotedblright{} and \emph{rogāre} \textquotedblleft{}to ask\textquotedblright{}.
\end{cousinweb}

\begin{grammarlens}[title={What you've built}]
One more word for telling, explaining and describing at B1.
\end{grammarlens}

\subsection*{Your turn}
\begin{itemize}
\raggedright
  \item \emph{Say it:} \emph{l'interrogazione}
  \item \emph{Say it:} \emph{l'interrogazione}, once more
  \item \emph{Say it:} a sentence of your own with \emph{l'interrogazione}
  \item \emph{Recall:} say the Italian for the exam session, then the Italian for to revise, then say \emph{l'interrogazione} again
\end{itemize}

\begin{tcolorbox}[breakable,colback=teal!4,colframe=teal!35!black,title={Before you move on}]
What does l'interrogazione mean? (\textquotedblleft{}The oral test (at school); the questioning\textquotedblright{}.) Say it once more.
\end{tcolorbox}
\section[memorizzare]{memorizzare — to memorise, to learn by heart}

\label{lesson:IT-C382-memorizzare}



\begin{quote}
Before the new one: say the Italian for to revise, then the Italian for the oral test.
\end{quote}

\subsection*{You'll want to know: memorizzare}
\textbf{memorizzare} — \textquotedblleft{}to memorise, to learn by heart\textquotedblright{}.

\begin{quote}
Devo memorizzare venti parole al giorno.
\end{quote}

\begin{quote}
\emph{I have to memorise twenty words a day.}
\end{quote}

On a phone, also to save: \emph{memorizzare un numero}. \emph{Imparare a memoria} is more colloquial.

\begin{cousinweb}
\emph{memorizzare} is built on \emph{memoria} \textquotedblleft{}memory\textquotedblright{}, a learned borrowing from Latin \emph{memoria} (Lewis \& Short), with the verb ending \emph{-izzare}; the modern word was probably helped along by English \emph{memorize}.
\end{cousinweb}

\begin{grammarlens}[title={What you've built}]
One more word for telling, explaining and describing at B1.
\end{grammarlens}

\subsection*{Your turn}
\begin{itemize}
\raggedright
  \item \emph{Say it:} \emph{memorizzare}
  \item \emph{Say it:} \emph{memorizzare}, once more
  \item \emph{Say it:} a sentence of your own with \emph{memorizzare}
  \item \emph{Recall:} say the Italian for to revise, then the Italian for the oral test, then say \emph{memorizzare} again
\end{itemize}

\begin{tcolorbox}[breakable,colback=teal!4,colframe=teal!35!black,title={Before you move on}]
What does memorizzare mean? (\textquotedblleft{}To memorise, to learn by heart\textquotedblright{}.) Say it once more.
\end{tcolorbox}
\section[gli appunti]{gli appunti — the notes (taken in class)}

\label{lesson:IT-C382-appunti}



\begin{quote}
Before the new one: say the Italian for the oral test, then the Italian for to memorise.
\end{quote}

\subsection*{You'll want to know: gli appunti}
\textbf{gli appunti} — \textquotedblleft{}the notes (taken in class)\textquotedblright{}.

\begin{quote}
Mi presti i tuoi appunti? Ieri ero malato.
\end{quote}

\begin{quote}
\emph{Can you lend me your notes? I was ill yesterday.}
\end{quote}

\emph{Prendere appunti} — to take notes. The singular \emph{appunto} also means a reproach; the adverb \emph{appunto} means exactly.

\begin{cousinweb}
\emph{appunti} is built on \emph{appuntare} \textquotedblleft{}to jot down, to pin\textquotedblright{}, from \emph{a-} and \emph{punto} \textquotedblleft{}point\textquotedblright{}, which is inherited from Latin \emph{punctum} \textquotedblleft{}a prick, a point\textquotedblright{} (Lewis \& Short).
\end{cousinweb}

\begin{grammarlens}[title={What you've built}]
One more word for telling, explaining and describing at B1.
\end{grammarlens}

\subsection*{Your turn}
\begin{itemize}
\raggedright
  \item \emph{Say it:} \emph{gli appunti}
  \item \emph{Say it:} \emph{gli appunti}, once more
  \item \emph{Say it:} a sentence of your own with \emph{gli appunti}
  \item \emph{Recall:} say the Italian for the oral test, then the Italian for to memorise, then say \emph{gli appunti} again
\end{itemize}

\begin{tcolorbox}[breakable,colback=teal!4,colframe=teal!35!black,title={Before you move on}]
What does gli appunti mean? (\textquotedblleft{}The notes (taken in class)\textquotedblright{}.) Say it once more.
\end{tcolorbox}
\section[il manuale]{il manuale — the textbook, the manual}

\label{lesson:IT-C382-manuale}



\begin{quote}
Before the new one: say the Italian for to memorise, then the Italian for the notes.
\end{quote}

\subsection*{You'll want to know: il manuale}
\textbf{il manuale} — \textquotedblleft{}the textbook, the manual\textquotedblright{}.

\begin{quote}
Il manuale di storia ha cinquecento pagine.
\end{quote}

\begin{quote}
\emph{The history textbook has five hundred pages.}
\end{quote}

As an adjective: \emph{lavoro manuale} — work with your hands.

\begin{cousinweb}
\emph{manuale} is a learned borrowing from Latin \emph{manuālis} \textquotedblleft{}of the hand, held in the hand\textquotedblright{} (Lewis \& Short), from \emph{manus} \textquotedblleft{}hand\textquotedblright{}: a book to keep at hand.
\end{cousinweb}

\begin{grammarlens}[title={What you've built}]
One more word for telling, explaining and describing at B1.
\end{grammarlens}

\subsection*{Your turn}
\begin{itemize}
\raggedright
  \item \emph{Say it:} \emph{il manuale}
  \item \emph{Say it:} \emph{il manuale}, once more
  \item \emph{Say it:} a sentence of your own with \emph{il manuale}
  \item \emph{Recall:} say the Italian for to memorise, then the Italian for the notes, then say \emph{il manuale} again
\end{itemize}

\begin{tcolorbox}[breakable,colback=teal!4,colframe=teal!35!black,title={Before you move on}]
What does il manuale mean? (\textquotedblleft{}The textbook, the manual\textquotedblright{}.) Say it once more.
\end{tcolorbox}
